% Lecture 3 notes: \input by notes.tex, \subimport'ed by the semester book.
% Exam questions (зачёт): 1 (деление с остатком больших чисел), 12 (алгоритм Евклида,
% расширенный алгоритм Евклида), 13 (обратный элемент кольца вычетов).
\begin{lecture}{3}{Деление с остатком и алгоритм Евклида}

\section{Зачем криптографии деление и алгоритм Евклида}

Почти любая криптографическая операция "--- это вычисление в кольце вычетов $\Zn{n}$: результат
каждого умножения больших чисел нужно \emph{привести по модулю} $n$, т.\,е. найти остаток от
деления на $n$. Поэтому деление с остатком "--- такая же базовая операция длинной арифметики, как
умножение (лекция~2), и её скорость напрямую определяет скорость шифрования.

Вторая базовая операция "--- \emph{обращение} по модулю. Например, в криптосистеме RSA
(семестр~2) открытый ключ "--- пара $(N, e)$, а секретная экспонента $d$ "--- это элемент, обратный
к $e$ по модулю $\varphi(N)$: $ed \equiv 1 \pmod{\varphi(N)}$. Для учебных параметров $p = 61$,
$q = 53$, $N = pq = 3233$, $\varphi(N) = 60 \cdot 52 = 3120$, $e = 17$ получаем $d = 2753$:
действительно, $17 \cdot 2753 = 46801 = 15 \cdot 3120 + 1$. Как найти $d$, не перебирая все
вычеты? Ответ даёт расширенный алгоритм Евклида "--- один из старейших алгоритмов (около 300 г.
до н.\,э.) и один из самых используемых: обращение по модулю нужно и в цифровых подписях
(DSA, ECDSA), и при сложении точек эллиптической кривой, и в методе
Монтгомери (лекция~5), и в китайской теореме об остатках.

План лекции: деление с остатком в $\ZZ$ и деление больших чисел (алгоритм~D Кнута), наибольший
общий делитель и алгоритм Евклида с оценкой числа шагов, расширенный алгоритм Евклида, обратимые
элементы кольца вычетов и теорема Эйлера, линейные сравнения.

\section{Деление с остатком}

\begin{theorem}[о делении с остатком]\label{l3:divrem}
Для любых целых $a$ и $b > 0$ существуют единственные целые $q$ и $r$ такие, что
\[
  a = qb + r, \qquad 0 \le r < b.
\]
При этом $q = \floor{a/b}$.
\end{theorem}

\begin{proof}
\emph{Существование.} Множество $S = \{a - kb : k \in \ZZ\} \cap \NN_0$ неотрицательных чисел вида
$a - kb$ непусто (при $k = -\abs{a}$ получаем $a + \abs{a} b \ge 0$). Пусть $r = a - qb$ "---
его наименьший элемент. Если бы $r \ge b$, то $r - b = a - (q + 1)b \in S$ и $r - b < r$ "---
противоречие. Итак, $0 \le r < b$.

\emph{Единственность.} Если $a = qb + r = q'b + r'$ и $0 \le r, r' < b$, то $b \mid r - r'$
и $\abs{r - r'} < b$, откуда $r = r'$ и $q = q'$.

Наконец, $0 \le r < b$ равносильно $q \le a/b < q + 1$, т.\,е. $q = \floor{a/b}$.
\end{proof}

Числа $q$ и $r$ называются \emph{неполным частным} и \emph{остатком}; пишут $r = a \bmod b$,
$q = a \operatorname{div} b$. Обратите внимание: остаток всегда неотрицателен, например
$-7 = (-3) \cdot 3 + 2$, т.\,е. $-7 \bmod 3 = 2$ (в языке Си оператор \texttt{\%} даёт для
отрицательного делимого отрицательный остаток $-1$, и это частый источник ошибок). Иногда
удобен \emph{абсолютно наименьший} остаток $-b/2 < r \le b/2$; он также существует и единствен.

\subsection{Деление на цифру}

Пусть $B$ "--- основание, $u = (u_{n-1} \dots u_0)_B$, $0 < v < B$ "--- однозначный делитель.
Деление «столбиком» на однозначное число выполняется одним проходом от старших разрядов.

\begin{algorithm}[деление на цифру]
\begin{algorithmic}[1]
\Require $u = (u_{n-1} \dots u_0)_B$, цифра $0 < v < B$.
\Ensure $q = (q_{n-1} \dots q_0)_B = \floor{u/v}$, $r = u \bmod v$.
\State $r \gets 0$
\For{$j = n - 1, n - 2, \dots, 0$}
  \State $t \gets rB + u_j$ \Comment{двухцифровое число}
  \State $q_j \gets \floor{t / v}$, \quad $r \gets t \bmod v$
\EndFor
\State \Return $q, r$
\end{algorithmic}
\end{algorithm}

Так как $r < v$, то $t = rB + u_j \le (v - 1)B + B - 1 < vB$, поэтому $q_j < B$ "--- это
действительно цифра, а $t$ "--- двухцифровое число, которое процессор делит на цифру одной
командой. Индукцией по $j$ проверяется инвариант
$(u_{n-1} \dots u_j)_B = (q_{n-1} \dots q_j)_B \cdot v + r$. Сложность "--- $\Theta(n)$.

\begin{example}
$9876543 : 7$: последовательные $t$ и $r$ равны
$9 \to (1, 2)$, $28 \to (4, 0)$, $7 \to (1, 0)$, $6 \to (0, 6)$, $65 \to (9, 2)$, $24 \to (3, 3)$,
$33 \to (4, 5)$ (в скобках "--- цифра частного и остаток), т.\,е. $9876543 = 1410934 \cdot 7 + 5$.
\end{example}

\section{Деление больших чисел: алгоритм D}

Пусть теперь $u = (u_{m+n-1} \dots u_0)_B$ и $v = (v_{n-1} \dots v_0)_B$, $v_{n-1} \ne 0$,
$n \ge 2$. Частное $\floor{u/v}$ имеет не более $m + 1$ цифр. Школьный метод находит их по одной,
начиная со старшей: на каждом шаге берётся \emph{окно} "--- $(n + 1)$-значное число $U$,
составленное из текущего остатка и очередной цифры делимого, и находится одна цифра
\[
  q = \floor{U / v},
\]
после чего $U - qv$ становится новым остатком. Если в начале шага остаток меньше $v$, то
$U < vB$ и $q \le B - 1$ "--- одна цифра. Вся трудность в том, как \emph{найти} $q$: в школе его
угадывают, а компьютеру нужна формула. Естественно оценить $q$ по старшим цифрам:
\begin{equation}\label{l3:qhat}
  \hat q = \min\left(\floor{\frac{U_n B + U_{n-1}}{v_{n-1}}},\ B - 1\right),
\end{equation}
где $U = (U_n U_{n-1} \dots U_0)_B$ (старшая цифра $U_n$ может быть нулём). Это деление
двухцифрового числа на цифру. Насколько $\hat q$ может отличаться от $q$?

\begin{theorem}[Кнут]\label{l3:qhat-thm}
Пусть $0 \le U < vB$ и $\hat q$ определено формулой~\eqref{l3:qhat}. Тогда
\begin{enumerate}
  \item $\hat q \ge q$;
  \item если $v_{n-1} \ge \floor{B/2}$, то $\hat q - 2 \le q \le \hat q$.
\end{enumerate}
\end{theorem}

\begin{proof}
1. Так как $q \le B - 1$, при $\hat q = B - 1$ утверждение очевидно. Иначе
$\hat q = \floor{(U_n B + U_{n-1})/v_{n-1}}$, и тогда $\hat q v_{n-1} \ge U_n B + U_{n-1} - v_{n-1} + 1$.
Отсюда
\[
\begin{aligned}
  U - \hat q v &\le U - \hat q v_{n-1} B^{n-1}\\
  &\le U_n B^n + U_{n-1} B^{n-1} + (B^{n-1} - 1) - (U_n B + U_{n-1} - v_{n-1} + 1) B^{n-1}\\
  &= v_{n-1} B^{n-1} - 1 < v.
\end{aligned}
\]
(Мы оценили младшие цифры окна: $U_{n-2} B^{n-2} + \cdots + U_0 \le B^{n-1} - 1$.) Итак,
$U - \hat q v < v$, т.\,е. $U/v < \hat q + 1$, и $q = \floor{U/v} \le \hat q$.

2. Предположим противное: $\hat q \ge q + 3$. Так как $v_{n-1} B^{n-1} > v - B^{n-1}$,
\[
  \hat q \le \frac{U_n B + U_{n-1}}{v_{n-1}} = \frac{U_n B^n + U_{n-1} B^{n-1}}{v_{n-1} B^{n-1}}
  \le \frac{U}{v_{n-1} B^{n-1}} < \frac{U}{v - B^{n-1}}.
\]
(Случай $v = B^{n-1}$ невозможен: тогда $q = \floor{U/B^{n-1}} = U_n B + U_{n-1} \ge \hat q$.)
Так как $q > U/v - 1$, получаем
\[
  3 \le \hat q - q < \frac{U}{v - B^{n-1}} - \Bigl(\frac{U}{v} - 1\Bigr)
  = \frac{U}{v} \cdot \frac{B^{n-1}}{v - B^{n-1}} + 1,
\]
откуда
\[
  \frac{U}{v} > 2\,\frac{v - B^{n-1}}{B^{n-1}} \ge 2(v_{n-1} - 1).
\]
Значит, $q = \floor{U/v} \ge 2(v_{n-1} - 1)$. С другой стороны, $q \le \hat q - 3 \le B - 4$.
Итак, $2v_{n-1} - 2 \le B - 4$, т.\,е. $v_{n-1} \le B/2 - 1 < \floor{B/2}$, что противоречит
условию.
\end{proof}

Итак, если старшая цифра делителя «большая», оценка $\hat q$ ошибается не более чем на $2$.
Условие $v_{n-1} \ge \floor{B/2}$ обеспечивается \emph{нормализацией}: делимое и делитель
умножаются на одну и ту же цифру $d$, от чего частное не меняется, а остаток умножается на $d$.

\begin{lemma}\label{l3:normalize}
Пусть $d = \floor{B/(v_{n-1} + 1)}$. Тогда $vd$ "--- по-прежнему $n$-значное число, и его старшая
цифра не меньше $\floor{B/2}$.
\end{lemma}

\begin{proof}
$vd < (v_{n-1} + 1) B^{n-1} d \le B \cdot B^{n-1} = B^n$, т.\,е. длина не увеличилась. Так как
$vd \ge v_{n-1} d B^{n-1}$, старшая цифра $vd$ не меньше $x d$, где $x = v_{n-1}$. Если
$x \ge \floor{B/2}$, то $d = 1$ и доказывать нечего. Пусть $x < \floor{B/2}$ и, от противного,
$xd \le \floor{B/2} - 1$. По определению $d$ имеем $(d + 1)(x + 1) > B$, т.\,е.
$xd \ge B - d - x$, откуда $d + x \ge B - xd \ge \ceil{B/2} + 1$. Но $d, x \ge 1$, поэтому
$(d - 1)(x - 1) \ge 0$, т.\,е. $xd \ge d + x - 1 \ge \ceil{B/2}$ "--- противоречие. Например, при $B = 10$: $v_{n-1} = 1, 2, 3, 4$ дают $d = 5, 3, 2, 2$ и старшие цифры
$vd$ не меньше $5, 6, 6, 8$.
\end{proof}

При $B = 2^{32}$ на практике берут $d = 2^s$: нормализация "--- сдвиг влево на число ведущих
нулевых битов $v_{n-1}$, так что старший бит делителя становится единицей.

\subsection{Уточнение оценки}

Кнут предложил дешёвую проверку, которая до вычитания отбрасывает почти все неверные
значения $\hat q$. Пусть $\hat r = U_n B + U_{n-1} - \hat q v_{n-1}$ "--- остаток от деления двух
старших цифр окна на $v_{n-1}$ (с учётом ограничения $\hat q \le B - 1$).

\begin{lemma}\label{l3:refine}
\begin{enumerate}
  \item Если $\hat q v_{n-2} > B \hat r + U_{n-2}$, то $q < \hat q$.
  \item Если $\hat q v_{n-2} \le B \hat r + U_{n-2}$, то $q \ge \hat q - 1$.
\end{enumerate}
\end{lemma}

\begin{proof}
Заметим, что $\hat q v_{n-1} + \hat r = U_n B + U_{n-1}$.

1. Отбрасывая младшие цифры $v$ и учитывая $U_{n-3} B^{n-3} + \cdots + U_0 < B^{n-2}$,
\[
  \hat q v \ge \hat q v_{n-1} B^{n-1} + \hat q v_{n-2} B^{n-2}
  \ge (U_n B + U_{n-1} - \hat r) B^{n-1} + (B \hat r + U_{n-2} + 1) B^{n-2}
  > U,
\]
т.\,е. $q < \hat q$.

2. При $\hat q = 0$ доказывать нечего. Пусть $\hat q \ge 1$; так как
$v < v_{n-1} B^{n-1} + (v_{n-2} + 1) B^{n-2}$,
\[
\begin{aligned}
  \hat q v &< \hat q v_{n-1} B^{n-1} + \hat q v_{n-2} B^{n-2} + \hat q B^{n-2}\\
  &\le (U_n B + U_{n-1} - \hat r) B^{n-1} + (B \hat r + U_{n-2}) B^{n-2} + \hat q B^{n-2}
  \le U + \hat q B^{n-2}.
\end{aligned}
\]
Так как $\hat q < B$, имеем $\hat q B^{n-2} < B^{n-1} \le v$, поэтому $U - (\hat q - 1) v > 0$
и $q \ge \hat q - 1$.
\end{proof}

Таким образом, проверку повторяют, уменьшая $\hat q$ на единицу (и увеличивая $\hat r$ на
$v_{n-1}$), пока она срабатывает; если $\hat r \ge B$, проверка заведомо не сработает
($\hat q v_{n-2} < B^2 \le B \hat r$), и её можно прекратить. После уточнения $\hat q \in \{q, q+1\}$,
а случай $\hat q = q + 1$ редок (при $B = 2^{32}$ "--- исчезающе редок) и обнаруживается при вычитании:
разность $U - \hat q v$ отрицательна. Тогда $\hat q$ уменьшают на единицу и прибавляют $v$ к
остатку («возврат»).

\begin{algorithm}[деление, алгоритм D Кнута]\label{l3:algD}
\begin{algorithmic}[1]
\Require $u = (u_{m+n-1} \dots u_0)_B$, $v = (v_{n-1} \dots v_0)_B$, $v_{n-1} \ne 0$, $n \ge 2$.
\Ensure $q = \floor{u/v} = (q_m \dots q_0)_B$, $r = u \bmod v$.
\State $d \gets \floor{B / (v_{n-1} + 1)}$; \ $u \gets ud = (u_{m+n} \dots u_0)_B$, \ $v \gets vd$
  \Comment{D1: нормализация}
\For{$j = m, m - 1, \dots, 0$}
  \State $\hat q \gets \min\bigl(\floor{(u_{j+n} B + u_{j+n-1}) / v_{n-1}},\ B - 1\bigr)$,
    \ $\hat r \gets u_{j+n} B + u_{j+n-1} - \hat q v_{n-1}$ \Comment{D3}
  \While{$\hat r < B$ и $\hat q v_{n-2} > B \hat r + u_{j+n-2}$}
    \State $\hat q \gets \hat q - 1$, \ $\hat r \gets \hat r + v_{n-1}$
  \EndWhile
  \State $(u_{j+n} \dots u_j)_B \gets (u_{j+n} \dots u_j)_B - \hat q v$ \Comment{D4: умножить и вычесть}
  \If{разность отрицательна}
    \State $\hat q \gets \hat q - 1$; \ $(u_{j+n} \dots u_j)_B \gets (u_{j+n} \dots u_j)_B + v$
      \Comment{D6: возврат}
  \EndIf
  \State $q_j \gets \hat q$
\EndFor
\State $r \gets (u_{n-1} \dots u_0)_B / d$ \Comment{D8: денормализация (деление на цифру)}
\State \Return $q, r$
\end{algorithmic}
\end{algorithm}

\begin{theorem}\label{l3:algD-cost}
Алгоритм D правильно вычисляет частное и остаток и выполняет $\bigO(mn)$ операций над цифрами
(при $m \ge 1$), т.\,е. $\bigO(N^2)$ для делимого длины $2N$ и делителя длины $N$.
\end{theorem}

\begin{proof}
\emph{Правильность.} Индукцией по шагам: перед шагом $j$ текущий остаток, т.\,е. число
$(u_{j+n} \dots u_{j+1})_B$, меньше делителя (вначале
$(u_{m+n} \dots u_{m+1})_B = \floor{ud / B^{m+1}}$,
а $ud / B^{m+1} < B^{n-1} d \le vd$, так как $u < B^{m+n}$ и $v \ge B^{n-1}$),
поэтому окно $U = (u_{j+n} \dots u_j)_B < vB$ и $q_j = \floor{U/v}$ "--- цифра. По
теореме~\ref{l3:qhat-thm} и лемме~\ref{l3:refine} после цикла уточнения $\hat q \in \{q_j, q_j + 1\}$,
а шаг возврата исправляет $\hat q = q_j + 1$. Новый остаток $U - q_j v$ меньше $v$. Как и при
делении на цифру, после шага $j = 0$ получаем $ud = qvd + r'$, $0 \le r' < vd$; значит, $q$ "---
частное от деления $u$ на $v$, а $r' = rd$.

\emph{Сложность.} Нормализация и денормализация стоят $\bigO(m + n)$. На каждом из $m + 1$ шагов
оценка $\hat q$ и не более двух уточнений стоят $\bigO(1)$, а умножение $\hat q v$, вычитание и
возможный возврат "--- $\bigO(n)$. Итого $\bigO(mn)$.
\end{proof}

Таким образом, школьные алгоритмы сложения, умножения и деления имеют сложность $\Theta(N)$,
$\Theta(N^2)$ и $\Theta(N^2)$ соответственно: деление $2N$-значного числа на $N$-значное
требует около $N^2$ умножений цифр "--- столько же, сколько умножение двух $N$-значных чисел.

\begin{example}\label{l3:algD-hand}
$u = 195264$, $v = 105$, $B = 10$, $n = 3$, $m = 3$. Нормализация: $d = \floor{10/2} = 5$,
$u' = 976320$ (с ведущим нулём $u'_6 = 0$), $v' = 525$, $v'_2 = 5$, $v'_1 = 2$.
\[
\begin{array}{c|c|l|r|c|c}
  j & U & \hat q & \hat q v' & q_j & U - q_j v' \\ \hline
  3 & 0976 & 1 & 525 & 1 & 451\\
  2 & 4513 & 9 \to 8 & 4200 & 8 & 313\\
  1 & 3132 & 6 & 3150 & 5^{*} & 507\\
  0 & 5070 & 10 \to 9 & 4725 & 9 & 345
\end{array}
\]
При $j = 2$: $\floor{45/5} = 9$, $\hat r = 0$, проверка $9 \cdot 2 = 18 > 10 \cdot 0 + 1$ срабатывает,
$\hat q = 8$, $\hat r = 5$, и $8 \cdot 2 = 16 \le 51$. При $j = 1$: $\hat q = 6$, $\hat r = 1$,
$6 \cdot 2 = 12 \le 13$ "--- проверка не срабатывает, но $6 \cdot 525 = 3150 > 3132$: нужен возврат
($^{*}$), $q_1 = 5$, $3132 - 2625 = 507$. При $j = 0$: $\floor{50/5} = 10 \ge B$, поэтому
$\hat q = 9$. Ответ: $q = 1859$, $r = 345/5 = 69$; проверка: $1859 \cdot 105 + 69 = 195264$.
\end{example}

\taskexample{division-knuth-d}

\checkyourself{Почему без нормализации оценка $\hat q$ может ошибаться сильно? Найдите при $B = 10$
окно $U$ и делитель $v$ со старшей цифрой $1$, для которых $\hat q - q \ge 5$.}

\checkyourself{Можно ли вместо $d = \floor{B/(v_{n-1}+1)}$ взять $d = \ceil{B/(2v_{n-1})}$?
Что может произойти с длиной делителя?}

\subsection{Деление через умножение}

Деление можно свести к умножению, и тогда оно наследует асимптотику быстрых алгоритмов
умножения. Обозначим через $M(N)$ сложность умножения $N$-значных чисел.

\emph{Метод Ньютона.} Чтобы вычислить $1/v$, применим метод Ньютона к функции
$f(x) = 1/x - v$:
\[
  x_{k+1} = x_k (2 - v x_k), \qquad 1 - v x_{k+1} = (1 - v x_k)^2.
\]
Погрешность возводится в квадрат, т.\,е. число верных цифр на каждом шаге удваивается.
Поскольку на $k$-м шаге достаточно работать с точностью около $2^k$ цифр, общая стоимость
вычисления $N$ цифр $1/v$ равна $\bigO(M(N) + M(N/2) + M(N/4) + \dots) = \bigO(M(N))$, после чего
$\floor{u/v}$ получается одним умножением и исправлением на $\pm 1$
(подробнее "--- \cite[разд.~3.4]{brent-zimmermann}, \cite[разд.~9.2.2]{crandall-pomerance}).

\emph{Редукция Барретта.} Если нужно много раз делить на \emph{один и тот же} модуль $n$ (как при
возведении в степень по модулю), обратная величина вычисляется один раз.

\begin{proposition}[Барретт]\label{l3:barrett}
Пусть $n$ имеет $k$ цифр по основанию $B$, $\mu = \floor{B^{2k}/n}$, $0 \le x < B^{2k}$, $q = \floor{x/n}$ и
\[
  \tilde q = \floor{\floor{x / B^{k-1}} \, \mu \, / \, B^{k+1}}.
\]
Тогда $q - 2 \le \tilde q \le q$, т.\,е. $x \bmod n = x - \tilde q n - \varepsilon n$,
$\varepsilon \in \{0, 1, 2\}$.
\end{proposition}

\begin{proof}
Пусть $\alpha = x/B^{k-1}$, $\beta = B^{2k}/n$, тогда $x/n = \alpha\beta / B^{k+1}$ и, так как
$\floor\alpha \floor\beta \le \alpha\beta$, получаем $\tilde q \le q$. С другой стороны,
$\alpha < B^{k+1}$ и $\beta \le B^{k+1}$ (так как $n \ge B^{k-1}$), поэтому
\[
  \frac{\floor\alpha \floor\beta}{B^{k+1}} > \frac{(\alpha - 1)(\beta - 1)}{B^{k+1}}
  > \frac{\alpha\beta}{B^{k+1}} - \frac{\alpha + \beta}{B^{k+1}} > \frac{x}{n} - 2,
\]
откуда $\tilde q \ge \floor{x/n} - 2$.
\end{proof}

Деления на степени $B$ "--- сдвиги, поэтому редукция стоит два умножения $k$-значных чисел
(точнее, нужны лишь старшие половины первого произведения и младшие "--- второго) и не более двух
вычитаний. Другой способ избавиться от деления "--- метод Монтгомери (лекция~5).

\section{Наибольший общий делитель и алгоритм Евклида}

Говорят, что $b$ \emph{делит} $a$ ($b \mid a$), если $a = bc$ для некоторого целого $c$.
\emph{Наибольшим общим делителем} $\gcd(a, b)$ целых чисел $a$, $b$, не равных нулю
одновременно, называется наибольшее целое, делящее и $a$, и $b$; по определению
$\gcd(0, 0) = 0$. Числа \emph{взаимно просты}, если $\gcd(a, b) = 1$.

Свойства, непосредственно следующие из определения: $\gcd(a, b) = \gcd(b, a) = \gcd(\abs a, \abs b)$,
$\gcd(a, 0) = \abs a$, и если $d \mid a$ и $d \mid b$, то $d \mid ax + by$ при любых целых $x$, $y$.

\begin{lemma}\label{l3:gcd-step}
Для любых целых $a$, $b$, $k$: $\gcd(a, b) = \gcd(a - kb, b)$. В частности, при $b > 0$
$\gcd(a, b) = \gcd(b, a \bmod b)$.
\end{lemma}

\begin{proof}
Общие делители пар $(a, b)$ и $(a - kb, b)$ совпадают: если $c \mid a$ и $c \mid b$, то
$c \mid a - kb$; обратно, если $c \mid a - kb$ и $c \mid b$, то $c \mid (a - kb) + kb = a$.
Значит, совпадают и наибольшие из них.
\end{proof}

\begin{algorithm}[Евклид]
\begin{algorithmic}[1]
\Require целые $a \ge b \ge 0$.
\Ensure $\gcd(a, b)$.
\While{$b \ne 0$}
  \State $(a, b) \gets (b,\ a \bmod b)$
\EndWhile
\State \Return $a$
\end{algorithmic}
\end{algorithm}

Положим $r_0 = a$, $r_1 = b$ и $r_{i+1} = r_{i-1} \bmod r_i$, $q_i = \floor{r_{i-1}/r_i}$:
\[
  r_{i-1} = q_i r_i + r_{i+1}, \qquad 0 \le r_{i+1} < r_i, \qquad i = 1, \dots, k,
\]
где $r_{k+1} = 0$ "--- первый нулевой остаток, а $k$ "--- число делений.

\begin{theorem}\label{l3:euclid}
Алгоритм Евклида останавливается, и $\gcd(a, b) = r_k$ "--- последний ненулевой остаток.
\end{theorem}

\begin{proof}
Остатки образуют строго убывающую последовательность неотрицательных целых
$r_1 > r_2 > \dots$, поэтому через конечное число шагов появится $r_{k+1} = 0$. По
лемме~\ref{l3:gcd-step} $\gcd(r_{i-1}, r_i) = \gcd(r_i, r_{i+1})$ на каждом шаге, так что
$\gcd(a, b) = \gcd(r_k, r_{k+1}) = \gcd(r_k, 0) = r_k$.
\end{proof}

\begin{example}
$a = 1071$, $b = 462$:
\[
  \euclidchain{1071}{462}
\]
и $\gcd(1071, 462) = 21$.
\end{example}

\subsection{Число шагов: теорема Ламе}

Насколько долго может работать алгоритм Евклида? Пусть $F_1 = F_2 = 1$, $F_{j+2} = F_{j+1} + F_j$ "---
числа Фибоначчи: $1, 1, 2, 3, 5, 8, 13, 21, 34, 55, 89, \dots$

\begin{theorem}[Ламе, 1844]\label{l3:lame}
Если $a > b > 0$ и алгоритм Евклида для $(a, b)$ выполняет $k$ делений, то $b \ge F_{k+1}$
и $a \ge F_{k+2}$.
\end{theorem}

\begin{proof}
Индукция по $k$. При $k = 1$: $b \ge 1 = F_2$ и $a > b$, т.\,е. $a \ge 2 = F_3$. Пусть $k \ge 2$ и
$a = q_1 b + r_2$. Так как делений больше одного, $r_2 > 0$, и для пары $(b, r_2)$, $b > r_2 > 0$,
алгоритм выполняет $k - 1$ делений. По предположению индукции $b \ge F_{k+1}$ и $r_2 \ge F_k$.
Поскольку $q_1 \ge 1$, получаем $a = q_1 b + r_2 \ge b + r_2 \ge F_{k+1} + F_k = F_{k+2}$.
\end{proof}

Оценка точна: для $a = F_{k+2}$, $b = F_{k+1}$ все частные равны $1$ (кроме последнего, равного $2$),
и делений ровно $k$. Например, для $(89, 55) = (F_{11}, F_{10})$ остатки
$34, 21, 13, 8, 5, 3, 2, 1, 0$ "--- девять делений.

\begin{corollary}\label{l3:lame-cor}
Пусть $\varphi = \frac{1 + \sqrt 5}{2} \approx 1{,}618$. Для $a > b > 0$ число делений $k$
удовлетворяет $k \le 1 + \log_\varphi b \approx 1 + 1{,}44 \log_2 b$. Если $b$ имеет $N$ десятичных
цифр, то $k \le 5N$.
\end{corollary}

\begin{proof}
Индукцией $F_j \ge \varphi^{j-2}$: $F_1 = 1 > \varphi^{-1}$, $F_2 = 1 = \varphi^0$, и
$F_{j+2} = F_{j+1} + F_j \ge \varphi^{j-1} + \varphi^{j-2} = \varphi^{j-2}(\varphi + 1) = \varphi^j$,
так как $\varphi^2 = \varphi + 1$. По теореме Ламе $b \ge F_{k+1} \ge \varphi^{k-1}$, откуда
первое неравенство ($\log_\varphi b = \log_2 b / \log_2 \varphi$, $1/\log_2\varphi \approx 1{,}44$).

Для второго заметим, что $F_{j+5} = 5F_{j+1} + 3F_j \ge 10 F_j$ при $j \ge 2$, так как
$F_{j+1} \ge \frac32 F_j$ при $j \ge 2$ (проверьте по индукции). Значит,
$F_{5N + 2} \ge 10^N F_2 = 10^N$. Если бы $k \ge 5N + 1$, то $b \ge F_{k+1} \ge F_{5N+2} \ge 10^N$,
а $b < 10^N$.
\end{proof}

\begin{theorem}\label{l3:euclid-cost}
Алгоритм Евклида для $a \ge b > 0$ выполняет $\bigO(\log b)$ делений и $\bigO(\log a \cdot \log b)$
битовых операций.
\end{theorem}

\begin{proof}
Число делений оценено в следствии~\ref{l3:lame-cor}. Деление $r_{i-1}$ на $r_i$ с частным $q_i$
школьным методом стоит $\bigO\bigl(\log r_i \cdot (1 + \log q_i)\bigr)$ битовых операций
(длина частного умножить на длину делителя). Так как $r_{i-1} \ge q_i r_i$, произведение всех
частных не превосходит $r_0 / r_k \le a$, поэтому
\[
  \sum_{i=1}^{k} \log r_i \,(1 + \log q_i) \le \log b \Bigl(k + \log \prod_i q_i\Bigr)
  = \bigO\bigl(\log b\,(\log b + \log a)\bigr) = \bigO(\log a \log b). \qedhere
\]
\end{proof}

Замечательно, что $\gcd$ двух $N$-битовых чисел вычисляется за $\bigO(N^2)$ "--- столько же,
сколько стоит \emph{одно} школьное умножение.

\checkyourself{Почему в доказательстве теоремы Ламе нужно условие $a > b$? Сколько делений
выполнит алгоритм для пары $(b, a)$, где $a > b$?}

\subsection{Бинарный алгоритм (Стейн)}

Деление больших чисел дороже вычитания и сдвига. Бинарный алгоритм Евклида (Дж.~Стейн, 1967)
использует только их, опираясь на свойства
\[
\begin{gathered}
  \gcd(2a, 2b) = 2\gcd(a, b), \qquad \gcd(2a, b) = \gcd(a, b) \ (b \text{ нечётно}),\\
  \gcd(a, b) = \gcd(\abs{a - b}, \min(a, b)),
\end{gathered}
\]
где в последнем случае для нечётных $a$, $b$ разность чётна и сразу делится на~$2$.

\begin{algorithm}[бинарный алгоритм Евклида]
\begin{algorithmic}[1]
\Require целые $a, b > 0$.
\Ensure $\gcd(a, b)$.
\State $s \gets$ наибольшее $s$ с $2^s \mid a$ и $2^s \mid b$; \ $a \gets a / 2^s$, $b \gets b / 2^s$
\State удалить из $a$ все множители $2$
\While{$b \ne 0$}
  \State удалить из $b$ все множители $2$ \Comment{теперь $a$, $b$ нечётны}
  \If{$a > b$} \State $(a, b) \gets (b, a)$ \EndIf
  \State $b \gets b - a$ \Comment{чётно}
\EndWhile
\State \Return $2^s a$
\end{algorithmic}
\end{algorithm}

На каждой итерации большее из двух нечётных чисел заменяется на число, не превосходящее
половины их разности, поэтому произведение $ab$ уменьшается по крайней мере вдвое, итераций
не больше $\log_2(ab)$, и каждая стоит $\bigO(\log a + \log b)$ битовых операций (сдвиг и вычитание):
сложность $\bigO(\log^2 \max(a, b))$, как и у классического алгоритма, но без делений.

\begin{example}
$\gcd(1071, 462)$: $462 = 2 \cdot 231$, $s = 0$; пары нечётных чисел
$(1071, 231) \to 840 = 2^3 \cdot 105$, $(231, 105) \to 126 = 2 \cdot 63$, $(105, 63) \to 42 = 2 \cdot 21$,
$(63, 21) \to 42 = 2 \cdot 21$, $(21, 21) \to 0$. Ответ $21$.
\end{example}

\section{Расширенный алгоритм Евклида}

Будем вместе с остатками $r_i$ хранить их запись в виде целочисленной линейной комбинации исходных
чисел:
\begin{equation}\label{l3:invariant}
  r_i = u_i a + v_i b.
\end{equation}
Для $r_0 = a$ и $r_1 = b$ годятся $(u_0, v_0) = (1, 0)$ и $(u_1, v_1) = (0, 1)$. Если
\eqref{l3:invariant} верно для $i - 1$ и $i$, то
\[
  r_{i+1} = r_{i-1} - q_i r_i = (u_{i-1} - q_i u_i)\,a + (v_{i-1} - q_i v_i)\,b,
\]
т.\,е. достаточно вычислять коэффициенты по \emph{тем же} рекуррентным формулам, что и остатки:
\begin{equation}\label{l3:xgcd-rec}
  u_{i+1} = u_{i-1} - q_i u_i, \qquad v_{i+1} = v_{i-1} - q_i v_i.
\end{equation}
Это и есть \emph{расширенный алгоритм Евклида}: в конце $r_k = \gcd(a, b) = u_k a + v_k b$.

\begin{algorithm}[расширенный алгоритм Евклида]
\begin{algorithmic}[1]
\Require целые $a \ge 0$, $b \ge 0$, не равные нулю одновременно.
\Ensure $d = \gcd(a, b)$ и $u, v \in \ZZ$ с $ua + vb = d$.
\State $(r_0, u_0, v_0) \gets (a, 1, 0)$, \ $(r_1, u_1, v_1) \gets (b, 0, 1)$
\While{$r_1 \ne 0$}
  \State $q \gets \floor{r_0 / r_1}$
  \State $(r_0, u_0, v_0,\ r_1, u_1, v_1) \gets (r_1, u_1, v_1,\ r_0 - q r_1, u_0 - q u_1, v_0 - q v_1)$
\EndWhile
\State \Return $(r_0, u_0, v_0)$
\end{algorithmic}
\end{algorithm}

Вычисления удобно записывать в таблицу: каждая новая строка получается из двух предыдущих
вычитанием предпоследней, умноженной на $q_i$. Для $a = 240$, $b = 46$:
\[
  \xgcdtable{240}{46} \qquad r_i = u_i \cdot 240 + v_i \cdot 46.
\]
Итак, $\gcd(240, 46) = 2 = -9 \cdot 240 + 47 \cdot 46$ (проверка: $-2160 + 2162 = 2$).

\begin{remark}
Если $a < b$, то $q_1 = 0$ и первый шаг просто меняет числа местами. Если нужен только
коэффициент $v$ (как при обращении по модулю $a$, см. ниже), столбец $u_i$ можно не вычислять:
он восстанавливается из $u = (d - vb)/a$.
\end{remark}

\begin{theorem}[соотношение Безу]\label{l3:bezout}
Для любых целых $a$, $b$ существуют целые $u$, $v$ такие, что $ua + vb = \gcd(a, b)$. Число
$\gcd(a, b)$ "--- наименьшее натуральное число вида $ua + vb$, и каждый общий делитель $a$ и $b$
делит $\gcd(a, b)$.
\end{theorem}

\begin{proof}
Существование $u$, $v$ даёт расширенный алгоритм Евклида (для отрицательных $a$, $b$ меняем
знаки коэффициентов). Любое число вида $xa + yb$ делится на $d = \gcd(a, b)$, поэтому положительные
такие числа не меньше $d$, а само $d$ имеет этот вид. Наконец, если $c \mid a$ и $c \mid b$, то
$c \mid ua + vb = d$.
\end{proof}

\begin{corollary}\label{l3:coprime-cor}
\begin{enumerate}
  \item $\gcd(a, b) = 1$ тогда и только тогда, когда $ua + vb = 1$ для некоторых целых $u$, $v$.
  \item (Лемма Евклида.) Если $a \mid bc$ и $\gcd(a, b) = 1$, то $a \mid c$.
  \item Если $d = \gcd(a, b)$, то $\gcd(a/d, b/d) = 1$.
\end{enumerate}
\end{corollary}

\begin{proof}
1. «Только тогда» "--- теорема~\ref{l3:bezout}; обратно, общий делитель $a$ и $b$ делит $1$.
2. $c = c(ua + vb) = a(cu) + (bc)v$ делится на $a$.
3. Разделим $ua + vb = d$ на $d$ и применим п.~1.
\end{proof}

Коэффициенты Безу определены неоднозначно: если $ua + vb = d$, то и
$(u + kb/d)\,a + (v - ka/d)\,b = d$ при любом целом $k$; по следствию~\ref{l3:coprime-cor} других
пар нет (если $ua + vb = u'a + v'b$, то $\frac{b}{d} \mid (u - u')\frac{a}{d}$, т.\,е.
$\frac{b}{d} \mid u - u'$). Расширенный алгоритм Евклида выдаёт пару с малыми коэффициентами.

\begin{theorem}\label{l3:xgcd-bounds}
Пусть $a > b > 0$, $d = \gcd(a, b)$, и $r_i, u_i, v_i$ "--- строки расширенного алгоритма Евклида,
$r_k = d$, $r_{k+1} = 0$. Тогда
\begin{enumerate}
  \item знаки чередуются: $u_i u_{i+1} \le 0$, $v_i v_{i+1} \le 0$, и
    $r_i \abs{u_{i+1}} + r_{i+1} \abs{u_i} = b$, \ $r_i \abs{v_{i+1}} + r_{i+1} \abs{v_i} = a$;
  \item $\abs{u_{k+1}} = b/d$, $\abs{v_{k+1}} = a/d$;
  \item $\abs{u_k} \le \dfrac{b}{2d}$, \ $\abs{v_k} \le \dfrac{a}{2d}$.
\end{enumerate}
\end{theorem}

\begin{proof}
1. Индукция по $i$. Для $i = 0$: $u_0 u_1 = v_0 v_1 = 0$, $r_0 \cdot 0 + r_1 \cdot 1 = b$ и
$r_0 \cdot 1 + r_1 \cdot 0 = a$. По \eqref{l3:xgcd-rec} $u_{i+2} = u_i - q_{i+1} u_{i+1}$, а числа $u_i$
и $-q_{i+1}u_{i+1}$ не имеют противоположных знаков (так как $u_i u_{i+1} \le 0$); поэтому
$\abs{u_{i+2}} = \abs{u_i} + q_{i+1}\abs{u_{i+1}}$ и $u_{i+1} u_{i+2} \le 0$. Тогда
\[
  r_{i+1}\abs{u_{i+2}} + r_{i+2}\abs{u_{i+1}}
  = r_{i+1}\abs{u_i} + q_{i+1} r_{i+1} \abs{u_{i+1}} + (r_i - q_{i+1} r_{i+1}) \abs{u_{i+1}}
  = r_{i+1}\abs{u_i} + r_i\abs{u_{i+1}} = b,
\]
и аналогично для $v$.

2. Подставим $i = k$: $r_k \abs{u_{k+1}} + 0 = b$, т.\,е. $\abs{u_{k+1}} = b/d$; так же для $v$.

3. Подставим $i = k - 1$: $r_{k-1}\abs{u_k} + d\abs{u_{k-1}} = b$, откуда $\abs{u_k} \le b / r_{k-1}$.
При $k \ge 2$ число $r_{k-1}$ кратно $d$ и больше $r_k = d$, т.\,е. $r_{k-1} \ge 2d$. При $k = 1$
имеем $b \mid a$, $u_1 = 0$, $v_1 = 1$, и $1 \le a/(2b)$, так как $a \ge 2b$. Для $v_k$ рассуждение
то же.
\end{proof}

Например, для $(240, 46)$: $\abs{-9} \le 46/4$, $\abs{47} \le 240/4$, а строка нулевого остатка
даёт $(u, v) = (23, -120) = (46/2, -240/2)$. В частности, коэффициенты по модулю не превосходят
исходных чисел и помещаются в машинное представление той же длины. Сложность расширенного
алгоритма "--- та же, $\bigO(\log a \log b)$ битовых операций: вычисление $u_i, v_i$ по формулам
\eqref{l3:xgcd-rec} стоит не больше, чем деления (оценка аналогична теореме~\ref{l3:euclid-cost}).

\taskexample{euclid-xgcd}

\checkyourself{Докажите, что $u_i v_{i+1} - u_{i+1} v_i = (-1)^i$. Почему отсюда следует, что
$\gcd(u_i, v_i) = 1$ для всех $i$?}

\checkyourself{Как изменить бинарный алгоритм Евклида, чтобы он вычислял и коэффициенты Безу?}

\section{Обратимые элементы кольца вычетов}

Целые $a$ и $b$ \emph{сравнимы по модулю} $n \ge 1$, $\congr{a}{b}{n}$, если $n \mid a - b$, т.\,е.
если они дают одинаковые остатки при делении на $n$. Сравнимость "--- отношение эквивалентности,
согласованное со сложением и умножением: если $a \equiv a'$, $b \equiv b'$, то
$a + b \equiv a' + b'$ и $ab \equiv a'b' \pmod n$ (так как $ab - a'b' = a(b - b') + b'(a - a')$).
Классы эквивалентности $[a] = a + n\ZZ$ "--- \emph{вычеты} по модулю $n$ "--- образуют
\emph{кольцо вычетов} $\Zn{n} = \{[0], [1], \dots, [n - 1]\}$; каждый класс представлен своим
остатком от деления на $n$ (теорема~\ref{l3:divrem}).

\begin{definition}
Элемент $[a] \in \Zn{n}$ называется \emph{обратимым}, если существует $[x]$ с
$\congr{ax}{1}{n}$; тогда $[x] = [a]^{-1}$ "--- \emph{обратный} (мультипликативно обратный)
элемент. Обратимые элементы образуют группу по умножению $\Znx{n}$ "--- \emph{группу обратимых
элементов}.
\end{definition}

Обратный элемент, если существует, единствен: из $ax \equiv 1 \equiv ax'$ следует
$x \equiv x(ax') = (xa)x' \equiv x'$. Произведение обратимых элементов обратимо:
$(ab)^{-1} = b^{-1}a^{-1}$.

\begin{theorem}\label{l3:invertible}
Вычет $[a] \in \Zn{n}$ обратим тогда и только тогда, когда $\gcd(a, n) = 1$. Если
$ua + vn = 1$, то $a^{-1} \equiv u \pmod n$. Если же $d = \gcd(a, n) > 1$, то $[a]$ "--- делитель
нуля: $a \cdot (n/d) \equiv 0 \pmod n$, $n/d \not\equiv 0$.
\end{theorem}

\begin{proof}
Если $\gcd(a, n) = 1$, то по теореме~\ref{l3:bezout} $ua + vn = 1$, т.\,е. $ua \equiv 1 \pmod n$.
Обратно, если $ax \equiv 1$, то $ax - 1 = kn$, т.\,е. $ax - kn = 1$, и всякий общий делитель
$a$ и $n$ делит $1$. Наконец, при $d > 1$: $a (n/d) = (a/d)\,n \equiv 0$, и $0 < n/d < n$.
(Делитель нуля не может быть обратим: из $ab \equiv 0$ и $xa \equiv 1$ следовало бы
$b \equiv xab \equiv 0$.)
\end{proof}

\begin{corollary}
$\Zn{n}$ "--- поле тогда и только тогда, когда $n = p$ простое. Поле из $p$ элементов
обозначается $\GF{p}$.
\end{corollary}

\begin{proof}
Все ненулевые вычеты $1, \dots, n - 1$ взаимно просты с $n$ в точности тогда, когда $n$ простое.
\end{proof}

\begin{algorithm}[обращение по модулю]\label{l3:inv-alg}
\begin{algorithmic}[1]
\Require $n \ge 2$, вычет $0 \le a < n$.
\Ensure $a^{-1} \bmod n$ или сообщение, что $a$ необратим.
\State расширенным алгоритмом Евклида найти $d = \gcd(n, a) = un + va$ \Comment{столбец $u$ не нужен}
\If{$d \ne 1$} \State \Return «необратим» \EndIf
\State \Return $v \bmod n$
\end{algorithmic}
\end{algorithm}

Сложность "--- $\bigO(\log^2 n)$ битовых операций; по теореме~\ref{l3:xgcd-bounds} $\abs v \le n/2$,
так что достаточно одного прибавления $n$, если $v < 0$.

\begin{example}\label{l3:rsa-inverse}
Секретная экспонента RSA из начала лекции: $d = 17^{-1} \bmod 3120$.
\[
  \xgcdtable{3120}{17} \qquad r_i = u_i \cdot 3120 + v_i \cdot 17.
\]
Строка последнего ненулевого остатка: $1 = 2 \cdot 3120 - 367 \cdot 17$, т.\,е.
$17^{-1} \equiv -367 \equiv 2753 \pmod{3120}$. Проверка: $17 \cdot 2753 = 46801 = 15 \cdot 3120 + 1$.
\end{example}

\taskexample{euclid-inverse}

\taskexample{euclid-noninvertible}

\subsection{Функция Эйлера и теорема Эйлера}

\begin{definition}
\emph{Функция Эйлера} $\varphi(n)$ "--- число вычетов $0 \le a < n$, взаимно простых с $n$, т.\,е.
$\varphi(n) = \abs{\Znx{n}}$.
\end{definition}

Например, $\varphi(p) = p - 1$ для простого $p$; $\varphi(p^e) = p^e - p^{e-1}$, так как среди
$0, \dots, p^e - 1$ ровно $p^{e-1}$ чисел делятся на $p$.

\begin{theorem}\label{l3:phi-mult}
Если $\gcd(m, n) = 1$, то $\varphi(mn) = \varphi(m)\varphi(n)$. Следовательно, для
$n = p_1^{e_1} \cdots p_s^{e_s}$
\[
  \varphi(n) = \prod_{i=1}^{s} p_i^{e_i - 1}(p_i - 1) = n \prod_{p \mid n}\Bigl(1 - \frac1p\Bigr).
\]
\end{theorem}

\begin{proof}
Рассмотрим отображение $f\colon \Zn{mn} \to \Zn{m} \times \Zn{n}$, $x \mapsto (x \bmod m, x \bmod n)$.
Оно инъективно: если $x \equiv y$ по модулю $m$ и по модулю $n$, то $m \mid x - y$ и
$n \mid x - y$, а так как $\gcd(m, n) = 1$, то по лемме Евклида $mn \mid x - y$. Мощности
множеств равны $mn$, поэтому $f$ "--- биекция. При этом $\gcd(x, mn) = 1$ тогда и только тогда,
когда $\gcd(x, m) = 1$ и $\gcd(x, n) = 1$ (простой делитель $mn$ делит $m$ или $n$), т.\,е. $f$
переводит обратимые вычеты в пары обратимых и наоборот. Отсюда
$\varphi(mn) = \varphi(m)\varphi(n)$. Формула для $\varphi(n)$ следует из мультипликативности и
значения на степенях простых. (Сюръективность $f$ "--- это китайская теорема об остатках,
семестр~2.)
\end{proof}

\begin{theorem}[Эйлер]\label{l3:euler}
Если $\gcd(a, n) = 1$, то $\congr{a^{\varphi(n)}}{1}{n}$.
\end{theorem}

\begin{proof}
Пусть $x_1, \dots, x_{\varphi(n)}$ "--- все обратимые вычеты по модулю $n$. Умножение на
обратимый $a$ "--- биекция $\Znx{n} \to \Znx{n}$ (обратное отображение "--- умножение на $a^{-1}$),
поэтому $ax_1, \dots, ax_{\varphi(n)}$ "--- те же вычеты в другом порядке, и
\[
  \prod_{i} (a x_i) \equiv \prod_i x_i \pmod n, \qquad\text{т.\,е.}\qquad
  a^{\varphi(n)} X \equiv X, \quad X = \prod_i x_i.
\]
Число $X$ обратимо как произведение обратимых; умножая на $X^{-1}$, получаем $a^{\varphi(n)} \equiv 1$.
\end{proof}

\begin{corollary}[малая теорема Ферма]
Если $p$ простое и $p \nmid a$, то $\congr{a^{p-1}}{1}{p}$; для любого $a$
$\congr{a^p}{a}{p}$.
\end{corollary}

Теорема Эйлера даёт второй способ обращения: при $\gcd(a, n) = 1$
\[
  a^{-1} \equiv a^{\varphi(n) - 1} \pmod n, \qquad\text{в частности}\qquad a^{-1} \equiv a^{p-2} \pmod p.
\]
Например, $\varphi(40) = 40 \cdot \frac12 \cdot \frac45 = 16$, и $7^{-1} \equiv 7^{15} \pmod{40}$:
$7^2 = 49 \equiv 9$, $7^4 \equiv 81 \equiv 1$, $7^{15} = (7^4)^3 \cdot 7^3 \equiv 343 \equiv 23$;
проверка: $7 \cdot 23 = 161 = 4 \cdot 40 + 1$. Однако этот способ хуже алгоритма Евклида:
возведение в степень требует $\bigO(\log n)$ умножений по модулю, т.\,е. $\bigO(\log^3 n)$ битовых
операций школьными методами (лекция~6), а для составного $n$ нужно ещё знать $\varphi(n)$, т.\,е.
фактически разложение $n$ на множители. Он полезен там, где важна \emph{регулярность}
вычислений (одинаковое время работы для всех $a$ в защищённых от атак по времени реализациях)
или где возведение в степень реализовано аппаратно.

\checkyourself{Как, зная $\varphi(N)$ для $N = pq$ ($p \ne q$ простые), найти $p$ и $q$? Почему
это означает, что вычисление $\varphi(N)$ не проще разложения $N$ на множители?}

\checkyourself{Найдите $\varphi(3120)$ ($3120 = 2^4 \cdot 3 \cdot 5 \cdot 13$). Сколько умножений
по модулю потребуется, чтобы найти $17^{-1} \bmod 3120$ по теореме Эйлера бинарным возведением
в степень? Сравните с числом шагов алгоритма Евклида в примере~\ref{l3:rsa-inverse}.}

\section{Линейные сравнения}

Обращение $a$ "--- это решение сравнения $ax \equiv 1 \pmod n$. Рассмотрим общее \emph{линейное
сравнение}
\begin{equation}\label{l3:lincongr}
  ax \equiv b \pmod n.
\end{equation}
Решения рассматриваются по модулю $n$: если $x$ "--- решение, то и $x + kn$.

\begin{theorem}\label{l3:linear}
Пусть $d = \gcd(a, n)$.
\begin{enumerate}
  \item Если $d \nmid b$, то сравнение~\eqref{l3:lincongr} не имеет решений.
  \item Если $d \mid b$, то оно имеет ровно $d$ решений по модулю $n$:
    \[
      x \equiv x_0 + k\frac{n}{d} \pmod n, \qquad k = 0, 1, \dots, d - 1,
    \]
    где $x_0$ "--- единственное решение сравнения
    $\frac{a}{d} x \equiv \frac{b}{d} \pmod{\frac{n}{d}}$. Если $ua + vn = d$, то
    $x_0 \equiv u \cdot \frac{b}{d} \pmod{\frac{n}{d}}$.
\end{enumerate}
\end{theorem}

\begin{proof}
1. Если $ax - kn = b$, то $d \mid b$, так как $d \mid a$ и $d \mid n$.

2. Положим $a' = a/d$, $b' = b/d$, $n' = n/d$; по следствию~\ref{l3:coprime-cor} $\gcd(a', n') = 1$.
Равенство $ax - b = kn$ равносильно $a'x - b' = kn'$, т.\,е. \eqref{l3:lincongr} равносильно
$a'x \equiv b' \pmod{n'}$. По теореме~\ref{l3:invertible} $a'$ обратим по модулю $n'$, и решения
этого сравнения "--- в точности числа $x \equiv x_0 = (a')^{-1} b' \pmod{n'}$. Из $ua + vn = d$
делением на $d$ получаем $ua' + vn' = 1$, т.\,е. $(a')^{-1} \equiv u \pmod {n'}$ и
$x_0 \equiv u b'$.

Остаётся сосчитать решения по модулю $n$. Числа $x_0 + kn'$, $k = 0, \dots, d - 1$, попарно
несравнимы по модулю $n$ (их разности по модулю меньше $n$ и не равны нулю), а любое
$x \equiv x_0 \pmod{n'}$ имеет вид $x_0 + jn'$ и сравнимо по модулю $n = dn'$ с
$x_0 + (j \bmod d) n'$.
\end{proof}

\begin{example}
$21x \equiv 14 \pmod{91}$: $\gcd(21, 91) = 7$ и $7 \mid 14$, поэтому решений семь. Сравнение
равносильно $3x \equiv 2 \pmod{13}$; $3^{-1} \equiv 9 \pmod{13}$ ($27 = 2 \cdot 13 + 1$),
$x_0 \equiv 9 \cdot 2 = 18 \equiv 5$. Решения: $x \in \{5, 18, 31, 44, 57, 70, 83\}$. А сравнение
$21x \equiv 15 \pmod{91}$ решений не имеет, так как $7 \nmid 15$.
\end{example}

\taskexample{euclid-linear-congruence}

\checkyourself{Сколько решений имеет сравнение $ax \equiv 0 \pmod n$? Опишите их через
$d = \gcd(a, n)$ и сравните с теоремой~\ref{l3:invertible}.}

\section{Что дальше}

\begin{itemize}
  \item Китайская теорема об остатках (семестр~2) обобщает биекцию из доказательства
    теоремы~\ref{l3:phi-mult}; её конструктивное доказательство "--- снова расширенный алгоритм
    Евклида.
  \item Метод Монтгомери (лекция~5) заменяет деление на модуль делением на степень основания;
    для него нужен обратный $-n^{-1} \bmod R$.
  \item Возведение в степень по модулю (лекция~6) "--- последовательность умножений и
    приведений по модулю; порядок элемента (лекция~7) уточняет теорему Эйлера.
  \item Существуют субквадратичные алгоритмы вычисления $\gcd$ (алгоритм Кнута"--~Шёнхаге,
    $\bigO(M(N)\log N)$), используемые в GMP для очень больших чисел.
\end{itemize}

\checkyourselfsection

\begin{furtherreading}
  \readitem[разд.~4.3.1]{knuth-taocp2}{алгоритм D, теоремы об оценке $\hat q$; разд.~4.5.2 "---
    алгоритм Евклида и бинарный алгоритм; разд.~4.5.3 "--- теорема Ламе и анализ числа шагов}
  \readitem[гл.~4, 10]{welschenbach}{деление с остатком и НОД, обратные элементы в реализации на Си}
  \readitem[разд.~1.2--1.4]{hps}{алгоритм Евклида, расширенный алгоритм, кольца вычетов, малая
    теорема Ферма}
  \readitem[разд.~2.1, 9.2]{crandall-pomerance}{алгоритм Евклида и обращение; редукция Барретта
    и деление методом Ньютона}
  \readitem[разд.~1.4, 1.6, 2.4]{brent-zimmermann}{деление, НОД, редукция Барретта}
  \readitem[разд.~14.2.5, 14.4]{hac}{алгоритмы деления, бинарного и расширенного алгоритма Евклида}
  \readitem{msu-ent}{алгоритм Евклида, сравнения, функция и теорема Эйлера (вики-курс МГУ)}
  \readitem{cp-algorithms}{статьи «Euclidean algorithm», «Extended Euclidean Algorithm»,
    «Modular Multiplicative Inverse», «Linear Congruence Equation» с реализациями}
\end{furtherreading}
\end{lecture}
